\documentclass[11pt,a4paper]{article}
\usepackage[margin=23mm]{geometry}
\usepackage{amsmath,amssymb,amsthm}
\setlength{\emergencystretch}{2em}
\newtheorem{theorem}{Theorem}
\newcommand{\gen}[1]{\langle#1\rangle}
\newcommand{\Alt}{\operatorname{Alt}}
\title{A two-point exterior also fails at odd-degree equality}
\author{Human Compression project: L044 continuation}
\date{}
\begin{document}\maketitle
Let $G=\gen{\rho_0,\ldots,\rho_{r-1}}\le\Alt(\Omega)$
be a transitive string C-group with connected diagram and all
maximal omissions intransitive. Let $D=\{1,\ldots,6\}$.
Assume the first three marks are an actual transitive head
supported on $D$, and that
\[
 \tau=\rho_3=(6\ b)(u\ v),\qquad
 R=G_{\ge3}=\Alt(\Omega\setminus\{1,\ldots,5\})
\]
are literal permutation statements, where $b,u,v$ are distinct
points outside $D$. Put $\sigma=\rho_4$, $H=G_{\ge5}$
and $K=G_{\ge6}$. Assume $\sigma$ fixes $D$ and $H$ fixes
$D\cup\{b\}$. No transitivity assumption is made on
$G_{\ge4}$ or on the support of $H$.

\begin{theorem}
If $r\ge6$, $n=|\Omega|\in\{2r,2r+1\}$, and every
colour $j\ge6$ has at least two full ambient-eligible edges,
then the displayed configuration is impossible.
\end{theorem}
\begin{proof}
Transitivity of $R$ forces $x=b^\sigma\ne b$, because
otherwise $R=\gen{\tau,\sigma,H}$ would preserve
$\{6,b\}$. The pure alternating overlap argument gives
\[
 \Alt(D\cup\{b\})\le G_{\le3},\qquad
 \Alt(D\cup\{b,x\})\le G_{\le4}.
\]
For completeness, conjugating $\tau$ by the transitive head
and multiplying two conjugates cancels $(u\ v)$, giving
all pure three-cycles through $b$ and two head points.
Conjugating the resulting alternating group by $\sigma$
and merging the two supports gives the second containment.
If $x^\tau\ne x$, one more such overlap puts a pure
$A_4$ on $\{6,b,x,x^\tau\}$ in
\[
                   G_{\le4}\cap R=\gen{\tau,\sigma},
\]
by full actual IP. A group generated by two involutions is
dihedral, and its odd-order elements commute; that pure
$A_4$ is impossible. Thus $\tau x=x$, so
$x\notin\{u,v\}$. The same intersection contains a
pure three-cycle $q$ on $\{6,b,x\}$.

The set $F=\{u,v\}$ is preserved by $H$: the latter fixes
$D,b$ and commutes with $\tau$. It cannot be preserved by
$\sigma$, since then all of $R$ would preserve it. If exactly
one point of $F$ moves outside, write $u^\sigma=w$ and
$v^\sigma=v$. The two invariant triples $\{6,b,x\}$
and $\{u,v,w\}$ each carry an order-three restriction of
$\tau\sigma$. Every element of odd order in its dihedral
group is a power of that rotation. A power trivial on the
second triple is divisible by three and trivial on the first.
This contradicts the pure cycle $q$.

Both points therefore move outside $F$. Write
$w=u^\sigma$, $z=v^\sigma$. The four points $u,v,w,z$
are distinct and disjoint from $D\cup\{b,x\}$.
The group $K$ fixes $D,b,x$, preserves $F$, and preserves
$F^\sigma=\{w,z\}$: it commutes with both $\tau$
and $\sigma$. Consequently it preserves
\[
 Y=\Omega\setminus(D\cup\{b,x,u,v,w,z\}),
                     \qquad |Y|=n-12.
\]
All ambient-eligible edges of colours $j\ge6$ lie in $Y$.
Indeed the only possible edges outside are on $F$ or
$F^\sigma$. The first are joined by $\tau\in G_j$;
the second by $\tau^\sigma=(6\ x)(w\ z)\in G_j$.
Neither is eligible for the full omission of $j$.

If $r>6$, apply the complete acquired localization theorem
to the terminal interval $I=\{6,\ldots,r-1\}$ and its
proper invariant set $Y$. It counts every actual $K$-orbit
in $Y$, including singletons, and gives
\[
                  2(r-6)\le |Y|-o_K(Y).
\]
For $n=2r$ this is impossible, since the left side is $|Y|$.
For $n=2r+1$, it forces $o_K(Y)=1$: the nonempty set $Y$
has odd size $2r-11$ and $K$ is transitive there.
The involution $\sigma$ preserves $Y$ and centralizes $K$.
An involution on an odd set has a fixed point, and a
centralizing permutation fixing one point of a transitive
action fixes every point. Hence $\sigma|_Y=1$.
If $r=6$, then $|Y|=0$ or $1$, and this conclusion holds
without the interval argument.

All points have now been accounted for. The permutation
$\sigma$ fixes $D$ and $Y$ and is exactly
\[
                         (b\ x)(u\ w)(v\ z).
\]
It is odd, contrary to $G\le\Alt(\Omega)$.
\end{proof}

\paragraph{Attribution and scope.}
The pure three-cycle and confined-edge argument extends the
accepted two-point exterior obstruction. Its full localization
proof remains paid. The new odd-degree step uses the forced
transitivity on $Y$ to eliminate the last possible odd factor
of $\sigma$. It uses no smaller-degree alternating rank
induction and no new computation. The literal natural $R$
and the exact two-point exterior remain hypotheses. This is
one case of the needed odd-degree separator refinement, not
a proof of that refinement or of the full L044 recovery.
No whole-original, formal or question credit is claimed.
\end{document}
